\documentclass[11pt]{article}
\usepackage[margin=1.1in]{geometry}
\usepackage[hidelinks]{hyperref}
\usepackage{parskip}
\pagestyle{empty}
\setlength{\parskip}{0.65\baselineskip}

\begin{document}

% ---- Letterhead: the only identity block in the letter ----
{\bfseries Aqi Dong}\\
Department of Management, Marketing and Operations\\
David B.\ O'Maley College of Business\\
Embry-Riddle Aeronautical University\\
Daytona Beach, FL 32114, USA\\
\texttt{donga2@erau.edu}

\vspace{1.2\baselineskip}
July 21, 2026

\vspace{1.2\baselineskip}
Editor-in-Chief\\
IEEE Transactions on Geoscience and Remote Sensing

\vspace{1.2\baselineskip}
Dear Editor-in-Chief,

I am pleased to submit my manuscript, \emph{``Tiled Dirichlet-Process Gaussian
Splatting for UAV Mapping with Calibrated Uncertainty,''} for consideration in
IEEE Transactions on Geoscience and Remote Sensing. The paper extends the
Dirichlet-process Gaussian-splatting model of my earlier methods paper
(arXiv:2607.10912) to survey-scale UAV mapping: independent per-tile
variational fits are fused by Hungarian matching under an exact conjugate
normal--inverse--Wishart marginal-likelihood score, merged in closed form, and
stitched by a partition-of-unity gate into a single \emph{normalized},
provenance-preserving scene-level predictive density. On three Hessigheim 3D
LiDAR epochs (59--129 million points) and a second 152-million-point
photogrammetric site, the stitched density matches the hard seam-deletion
rules of current tiled practice on held-out log-density while deleting
nothing, and the resulting calibrated uncertainty is then spent on four
mapping tasks: error sparsification, mesh-referenced accuracy stratification,
change detection with a model-derived level of detection, and re-flight
planning, all on CPU-only laptop hardware.

I believe TGRS is the right home for this work because its contribution is
operational rather than a methodological display: tiled processing of large UAV
surveys is standard practice in this community, the seam problem it addresses
is a practice problem, and every evaluation is a mapping-product evaluation
(held-out predictive density, empirical interval coverage, geometric
accuracy against the benchmark's reference mesh, null-calibrated change
detection on repeat LiDAR epochs, and acquisition planning) rather than a
rendering benchmark. The datasets (Hessigheim 3D, GauU-Scene, Mill 19,
STPLS3D) and the tasks are those of the remote-sensing literature the paper
engages.

I have tried to make the paper's claims exactly as strong as its evidence and
no stronger, and its controls and ablations are reported at parity: seam
deletion is empirically near-costless on these scenes and the paper says so;
seam-band parity and seam calibration hold at production capacity on the H3D
epochs but do not replicate at the second site, and the abstract carries that
caveat; the sparsification signal's point-density confound is decomposed
rather than hidden; the change-detection level of detection has a stated
${\approx}1$~m representation-limit floor; and the re-flight result is
dissected against oracle, random, and mass-only controls, including the
regime where the mass-only control matches the planner. A merge ablation and
a no-merge refit isolate what the exact conjugate merge contributes. Complete
code, fixed-seed reproduction commands, and machine-readable experiment
records accompany the submission.

This manuscript is not under consideration elsewhere; the methods paper it
builds on is cited and none of its figures or results are reused. I have no
conflicts of interest to declare.

Thank you for your consideration. I look forward to the reviewers' comments.

\vspace{1.2\baselineskip}
Sincerely,

\vspace{1.6\baselineskip}
Aqi Dong

\end{document}
